% Abstrak (Bahasa Indonesia)
\begin{center}
{\bfseries\large ABSTRAK}
\end{center}

\vspace{0.5cm}

\begin{spacing}{1}

\begin{center}
{\bfseries PENGEMBANGAN COMPANIONSHIP CHATBOT DENGAN ASESMEN SKRINING DEPRESI MELALUI PERCAKAPAN DAN MEMORI JANGKA PANJANG HIBRID BERBASIS KNOWLEDGE GRAPH UNTUK MAHASISWA INDONESIA}
\end{center}

\vspace{0.3cm}

\begin{center}
Oleh \\
Mohammad Nugraha Eka Prawira \\
NIM: 13522001
\end{center}

\vspace{0.5cm}

Mahasiswa Indonesia dapat menghadapi tekanan akademik dan emosional, tetapi dukungan profesional tidak selalu mudah diakses. \textit{Companionship chatbot} dapat menyediakan ruang pendampingan non-klinis yang mudah dijangkau apabila mampu menjaga percakapan tetap empatik, menawarkan skrining depresi secara hati-hati, dan mempertahankan konteks pengguna lintas sesi. Tugas akhir ini mengembangkan AXIS, sebuah \textit{companion chatbot} non-klinis yang dikalibrasi untuk konteks mahasiswa Indonesia. Sistem menggabungkan percakapan reflektif berprinsip \textit{Cognitive Behavioral Therapy}, pemeriksaan suasana hati harian dan asesmen skrining depresi PHQ-9 secara percakapan, serta memori jangka panjang hibrid dengan \textit{knowledge graph} sebagai struktur relasional utama dan pencarian semantik sebagai mekanisme penemuan kandidat. PHQ-9 diposisikan sebagai skrining, bukan diagnosis. Implementasi menggunakan orkestrasi percakapan berbasis graf keputusan dan batas keselamatan berlapis. Evaluasi dilakukan melalui validasi kontrak, \textit{benchmark} berlabel, penilaian buta \textit{LLM-as-judge}, dan ablasi terkendali. Hasil utama dikelompokkan ke dalam empat blok. Pada percakapan, 54 kontrak kebijakan dialog CBT lulus. Dalam 24 skenario penilaian buta, AXIS dipilih pada 14 penilaian, terutama pada kondisi bermemori, tetapi kesepakatan preferensi antarkonfigurasi penilai rendah (kappa $-0,04$), sehingga hasil ini tidak mendukung klaim keunggulan percakapan secara umum. Pada keselamatan, 80 kontrak lulus; sensitivitas deteksi pada 50 pesan mencapai 0,654 dan spesifisitas 0,917, serta seluruh respons setelah risiko terdeteksi memenuhi rubrik keselamatan. Pada PHQ-9, 130 kontrak lulus. Dari 80 jawaban bebas, 73 jawaban dengan frekuensi tegas dipetakan sesuai label acuan, sedangkan tujuh jawaban ambigu diarahkan ke klarifikasi. Pada memori, 47 kontrak lulus. Kondisi hibrid dan vektor-saja menghasilkan \textit{Recall@5} yang sama pada 15 kueri parafrase, tetapi kondisi hibrid berhasil pada satu kueri yang membutuhkan identitas relasi graf. Kalibrasi eksternal pada LongMemEval\_S menghasilkan \textit{Recall@5} 0,972 dan nDCG@5 0,855. Pengujian skenario-dan-hasil terhadap penulis memori produksi menghasilkan presisi 0,69, \textit{recall} 0,75, dan \textit{update correctness} 0,75, sedangkan penggunaan memori pada jawaban mencapai \textit{grounded-answer rate} 1,00, \textit{false-personalization rate} 0,00, dan \textit{stale-belief rate} 0,00 pada proksi berukuran kecil. Hasil tersebut menunjukkan kelayakan rekayasa purwarupa dan nilai tambah konteks graf pada kasus relasional tertentu, tetapi belum membuktikan efektivitas klinis, pengalaman pengguna mahasiswa Indonesia, atau keunggulan graf secara umum. Evaluasi pengguna nyata belum dilaksanakan dan menjadi pekerjaan lanjutan.

\vspace{0.5cm}

\noindent {\bfseries Kata Kunci:} \textit{companionship chatbot}, asesmen skrining depresi, PHQ-9, memori jangka panjang hibrid, \textit{knowledge graph}, mahasiswa Indonesia
\end{spacing}
